\chapter{多尺度食品热状态：从多孔介质到决策流形}
\label{ch:physics}

蒸烤失败的物理原因不是“温度没到”，而是四个场没有被同时看见：温度、水分、反应进度、结构。本章给出这些场的连续介质方程、它们如何耦合，以及何种低维状态足够支撑停机与探针决策。方程继承多孔介质体积平均与食品反应动力学的标准结构；原创点是把决策所需的泛函收成一条可估计流形，并标明哪些参数必须在线更新。

\section{计算域与状态}

食品区域记为 $\Omega_f\subset\R^3$，腔体气体区域记为 $\Omega_c$，交界面 $\Gamma=\partial\Omega_f$。时间 $t\in[0,t_f]$。食品视为可变孔隙的三相介质：固体基质、液态水、气体（干空气与水蒸气）。体积分数满足
\begin{equation}
\varepsilon_s+\varepsilon_l+\varepsilon_g=1,\qquad \varepsilon_s=\varepsilon_s(S)>0.
\end{equation}
$S(\vect{r},t)$ 是结构变量，综合收缩、孔隙闭合与表面结壳，不把它伪装成某一个分子机理。

\begin{definition}[全场状态]
\label{def:fullstate}
\begin{equation}
\label{eq:state}
\vect{X}(\vect{r},t)=\big(T,\,p_g,\,\varepsilon_l,\,\omega_v,\,\vect{\alpha},\,S\big),
\end{equation}
其中 $T$ 为温度，$p_g$ 为气相压力，$\varepsilon_l$ 为液态水体积分数，$\omega_v$ 为气相中水蒸气质量分数，反应进度
\begin{equation}
\vect{\alpha}=(\alpha_p,\alpha_g,\alpha_b,\alpha_m)\in[0,1]^4
\end{equation}
依次表示蛋白质变性、淀粉糊化、美拉德褐变与微生物致死当量的归一化进度。
\end{definition}

腔体状态不进入食品场，而作为边界输入 $\vect{u}_c=(T_\infty,Y_\infty,\vect{v}_\infty,\dot m_{\mathrm{st}},P_{\mathrm{heat}})$。

\section{宏观：热与三相质量}

采用 Whitaker 型体积平均。忽略热局部非平衡时，能量方程写为
\begin{equation}
\label{eq:energy}
(\rho c)_{\mathrm{eff}}\frac{\partial T}{\partial t}
=\nabla\!\cdot\!\big(k_{\mathrm{eff}}\nabla T\big)
-\sum_{i\in\{l,v,a\}}\nabla\!\cdot\!(h_i\vect{n}_i)
-\dot m_{\mathrm{phase}}\,\Delta h_{\mathrm{vap}}
+q_{\mathrm{reac}}.
\end{equation}
有效容积热容与导热系数为体积分数的加权，加权系数随 $\varepsilon_l$ 与 $S$ 变化，因此“同一种肉”在表面变干后导热下降。这是烤不匀与外干内生的机制，不是边界条件的附加说明。

液相与蒸汽：
\begin{align}
\label{eq:liquid}
\frac{\partial(\varepsilon_l\rho_l)}{\partial t}
&=-\nabla\!\cdot\!\vect{n}_l-\dot m_{\mathrm{phase}},\\
\label{eq:vapor}
\frac{\partial(\varepsilon_g\rho_g\omega_v)}{\partial t}
&=-\nabla\!\cdot\!\vect{n}_v+\dot m_{\mathrm{phase}}.
\end{align}
通量由 Darcy 渗流与 Fick 扩散合成：
\begin{align}
\vect{n}_l&=-\rho_l\frac{k_{rl}K}{\mu_l}\nabla p_l,\\
\vect{n}_v&=-\rho_g\omega_v\frac{k_{rg}K}{\mu_g}\nabla p_g
-\rho_g D_{\mathrm{eff}}\nabla\omega_v.
\end{align}
相变速率由局部偏离平衡驱动，而不是写成无限快的代数约束。无限快约束会使逆问题在数值上把任何温度误差直接折成水分误差。
\begin{equation}
\label{eq:phase}
\dot m_{\mathrm{phase}}
=K_{\mathrm{ev}}\,a_w(\varepsilon_l,T)\,
\big(p_{\mathrm{sat}}(T)-p_v\big)_+
-K_{\mathrm{cd}}\,(p_v-p_{\mathrm{sat}}(T))_+.
\end{equation}
$a_w$ 是水分活度。蒸制时腔体 $p_v$ 高，\eqref{eq:phase} 的冷凝项在表面释放潜热，这是蒸汽比干热更快把热量送进食品的原因。烤制时 $p_v$ 低，蒸发项吸热并留下干壳。

气相总压由干空气与蒸汽的状态方程闭合。干空气通量与 \eqref{eq:vapor} 对称：把 $\omega_v$ 换为干空气质量分数，相变源项改号。边界按温度、蒸汽质量分数、压力和辐射四类列出，见附录的闭合表。

\section{界面：腔体如何进入食品}

界面 $\Gamma$ 上的能量边界是对流、辐射与相变潜热之和：
\begin{equation}
\label{eq:bc}
-k_{\mathrm{eff}}\nabla T\!\cdot\!\vect{n}
=h\big(|\vect{v}_\infty|\big)\,(T-T_\infty)
+\epsilon\sigma(T^4-T_w^4)
+\dot m_{\Gamma}\Delta h_{\mathrm{vap}}.
\end{equation}
$h$ 随风速与冷凝状态切换：干对流用 Nusselt 关联，表面冷凝时改用膜状冷凝量级。两种关联的切换由腔体露点与表面温度的符号决定。蒸汽探针之所以携带信息，是因为它改变 \eqref{eq:bc} 的主项，而不是仅仅改变 $T_\infty$ 一个数。

质量边界：液态水不穿过 $\Gamma$，蒸汽通量与腔体湿度差成正比，冷凝液作为 $\varepsilon_l$ 的界面源。排汽阀打开时 $Y_\infty$ 下降，蒸发恢复。漏汽在模型里是 $Y_\infty$ 的偏差与门封热流，不单列成“故障文案”。

\section{介观与微观：反应进度不是温度的别名}

每一类反应满足
\begin{equation}
\label{eq:kinetics}
\frac{\partial\alpha_i}{\partial t}
=A_i\exp\!\Big(-\frac{E_{a,i}}{RT}\Big)
(1-\alpha_i)^{n_i}\,g_i(\varepsilon_l,S),
\end{equation}
$g_p$ 在自由水不足时下降（表面先干、蛋白凝固变慢或转成烘烤），$g_g$ 仅在淀粉与自由水同时存在时非零，$g_b$ 在低水分、高温时上升。因此同一条中心温度曲线可以对应不同的表面褐变。把“熟了”定义成 $T>T^\star$ 会同时弄错蒸与烤。

微生物安全不另造机理，使用冷点温度的等效时间，并归一化到 $\alpha_m$：
\begin{equation}
\label{eq:F}
\frac{\mathrm{d}F}{\mathrm{d}t}=10^{(T_{\mathrm{cold}}-T_{\mathrm{ref}})/z},\qquad
\alpha_m=1-\exp(-F/F_{\mathrm{req}}).
\end{equation}
$T_{\mathrm{cold}}(t)=\operatorname{ess\,inf}_{\vect{r}\in\Omega_f}T(\vect{r},t)$。冷点通常不在几何中心：蒸汽从上方冷凝、加热管在背部时，冷点会偏到背风侧。场模型必须允许冷点位置随 $\vect{u}_c$ 移动，不能写死在中心节点。

结构演化用经验但可辨识的 ODE，避免引入无法观测的细胞尺度变量：
\begin{equation}
\label{eq:structure}
\frac{\partial S}{\partial t}
=\kappa_s(1-S)\,\langle\alpha_p-\alpha_s^\star\rangle_+
+\kappa_c\,\langle M_{c}- \varepsilon_l\rangle_+,
\end{equation}
第一项是蛋白收缩，第二项是干燥结壳。$S$ 回调 $k_{\mathrm{eff}}$、$K$ 与 $h$。耦合是闭环的：干壳降低内部升温速度，使 $\alpha_m$ 落后于表面 $\alpha_b$。

\section{决策流形：估计什么，不估计什么}

全场 $\vect{X}$ 不能、也不需要实时重构到网格精度。停机、探针与解释只依赖一组泛函。

\begin{definition}[决策状态]
\label{def:decision}
\begin{equation}
\label{eq:z}
\vect{z}(t)=\big(
T_{\mathrm{cold}},\,
\bar T_s,\,
\sigma_T,\,
\bar\varepsilon_l,\,
\varepsilon_{l,s},\,
\alpha_m,\,
\alpha_{b,s},\,
S_s,\,
m_{\mathrm{loss}}
\big).
\end{equation}
依次为冷点温度、表面均温、温度场标准差、平均液相分数、表面液相分数、致死进度、表面褐变、表面结壳与累计质量损失。
\end{definition}

\begin{proposition}[决策充分性，工程表述]
若控制器的阶段代价与约束都可以写成 $\vect{z}$ 的 Lipschitz 函数，并且 $\vect{z}$ 的后验不确定度低于该 Lipschitz 常数所对应的阈值，则进一步细化网格场不会改变已投影的安全动作。
\end{proposition}

论证分为三步。动作是代价的 $\epsilon$-最优解。代价对 $\vect{z}$ 的变化由 Lipschitz 常数控制。后验期望代价之差小于动作切换的滞回带时，投影后的 $u$ 不变。滞回带属于决策层，安全不等式属于执行层，二者不共用阈值。因此增加红外像素而 $\vect{z}$ 的分位区间不变时，执行量保持不变。

全场到 $\vect{z}$ 的映射记为 $\Phi$。在线模型可以是：
\begin{enumerate}
\item 一维径向多孔介质，适合回转体食材的实时预测；
\item 两节点热容网络加一个集总水分仓，适合探针响应的秒级辨识；
\item 预计算参数化的神经算子，适合在食材类别内插值，见第 \ref{ch:evolve} 章。
\end{enumerate}
三者必须在 $\vect{z}$ 上可互相比对。不允许只在腔温上对齐。

\section{两节点模型：后续辨识的工作方程}

主动探针的可辨识性分析使用下面的最小闭合系统。表面节点 $1$、内部节点 $2$：
\begin{align}
\label{eq:twonode}
C_1\dot T_1
&=\frac{T_\infty-T_1}{R_{\mathrm{ext}}(v)}+\frac{T_2-T_1}{R_{\mathrm{int}}}
+q_{\mathrm{lat}}(\dot m_\Gamma),\\
C_2\dot T_2
&=\frac{T_1-T_2}{R_{\mathrm{int}}},\\
\dot m
&=-k_m\big(a_w(m)\,p_{\mathrm{sat}}(T_1)-p_\infty\big).
\end{align}
$R_{\mathrm{ext}}$ 由风速与是否冷凝切换，$R_{\mathrm{int}}$ 随结壳增大，$C_1,C_2$ 随含水率下降。$q_{\mathrm{lat}}$ 在冷凝时为正、蒸发时为负。

该模型不是腔体的数字孪生本体。它是信息探针的局部线性化载体：用来判断一次短脉冲能不能把 $R_{\mathrm{int}}/C_2$ 从 $R_{\mathrm{ext}}$ 里分开。全场方程 \eqref{eq:energy}--\eqref{eq:structure} 负责离线生成真值与训练算子。

\section{参数分层}

\begin{table}[H]
\centering
\caption{参数：哪些锁死，哪些在线更新}
\begin{tabularx}{\textwidth}{lXll}
\toprule
符号 & 含义 & 来源 & 在线 \\
\midrule
$p_{\mathrm{sat}}(T)$、$\Delta h_{\mathrm{vap}}$ & 水的物性 & 蒸汽表 & 否 \\
$A_i,E_{a,i},n_i$ & 反应动力学 & 类别先验 & 弱更新 \\
$k_{\mathrm{eff}},K,D_{\mathrm{eff}}$ & 有效传递 & 先验 + 本次 & 是 \\
$R_{\mathrm{ext}},h(v)$ & 腔体边界 & 探针 & 是 \\
$C_1,C_2,R_{\mathrm{int}}$ & 热容与内部热阻 & 探针 & 是 \\
$K_{\mathrm{ev}},K_{\mathrm{cd}}$ & 相变弛豫 & 探针 & 是 \\
几何厚度 $L$ & 冷点路径长度 & 视觉或用户 & 本次固定 \\
\bottomrule
\end{tabularx}
\end{table}

锁死蒸汽表，是为了不让网络用任意函数拟合潜热。在线只更新传递与边界。

\section{非均质与界面的闭合关系}

两节点模型用于可辨识性。真值模型在下列关系上比两节点更完整，这些关系不在在线参数里自由拟合。

水分活度采用 GAB 等温线，而不是把含水率直接当成蒸汽压：
\begin{equation}
\label{eq:gab}
M=\frac{M_0\,C\,K\,a_w}{(1-K a_w)\,(1-K a_w+C K a_w)}.
\end{equation}
$M_0,C,K$ 的类别区间由第 \ref{ch:data} 章的公开等温线约束；肉类在该数据中缺席，不借用淀粉参数冒充。温度移动按 Clausius--Clapeyron 修正吸附热，并标明外推超出测定温度的部分。表面蒸发驱动力使用 $a_w p_{\mathrm{sat}}(T_s)-p_v$，因此同一含水率在淀粉与蛋白上对应不同的界面蒸汽压。

有效导热写成各向异性张量，纤维或纹理方向 $\vect{e}$ 与法向分开：
\begin{equation}
\label{eq:aniso}
\vect{k}_{\mathrm{eff}}=k_\perp(M,S)\,\vect{I}+(k_\parallel-k_\perp)\,\vect{e}\otimes\vect{e}.
\end{equation}
脂肪夹层作为并联热阻 $R_{\mathrm{fat}}=L_{\mathrm{fat}}/(k_{\mathrm{fat}}A)$ 加到冷点路径上。冷点可以落在脂肪背风侧，而不必是几何中心。

收缩使计算域厚度随蛋白进度下降：
\begin{equation}
\label{eq:shrink}
\dot L=-\beta_s L_0\,\dot\alpha_p,\qquad K=K_0\Big(\frac{L}{L_0}\Big)^{n_K}\varepsilon_g^{n_\varepsilon}.
\end{equation}
渗透率随孔隙与收缩一起变，所以后期内部蒸汽更难排出，热阻上升不只有结壳一项。

腔壁辐射不用单一环境温度概括。加热管、食物表面与腔壁之间用一个角系数：
\begin{equation}
q_{\mathrm{rad}}=F_{h,s}\epsilon_s\sigma(T_h^4-T_s^4)+F_{w,s}\epsilon_s\sigma(T_w^4-T_s^4).
\end{equation}
$F_{h,s}$ 由管与料层的相对位置决定。热像观测中的反射项与 $F_{h,s}T_h^4$ 是同一几何量，标定一次，两处共用。

相变仍保持有限速率 \eqref{eq:phase}。无限快平衡会使温度残差被直接解释成水分残差，破坏探针的可分离性。

\section{三个原型的方程差异}

三个原型共用 \eqref{eq:energy}--\eqref{eq:structure}，只改变哪些项为主、哪些项关闭。这样 S1 的差异来自物理，而不是三套互不相关的公式。

\begin{table}[H]
\centering
\caption{原型开关}
\begin{tabularx}{\textwidth}{lXXX}
\toprule
项 & 肉糜板 & 土豆 & 面团 \\
\midrule
主反应 & $\alpha_p$ 与 $\alpha_m$ & $\alpha_g$，需自由水 & $\alpha_g$ 与后期 $\alpha_b$ \\
等温线 & 开放区间，无公开肉类曲线 & 第 \ref{ch:data} 章土豆清洁拟合的温度趋势 & 谷物清洁拟合的区间 \\
各向异性 & $k_\parallel/k_\perp$ 取文献区间 & 近似各向同性 & 气孔方向可有弱各向异性 \\
脂肪层 & T2 中可开启 & 关闭 & 关闭 \\
收缩 & \eqref{eq:shrink} 开启 & 弱 & 体积先胀后定，作为 $L(t)$ 的已知函数 \\
\bottomrule
\end{tabularx}
\end{table}

面团的膨胀不反解成内部状态。它作为几何输入，避免把发酵气体模型和热加工认知绑成一个估计问题。肉糜板的等温线在公开数据到达前保持宽区间，S1 报告该区间两端的策略排序是否一致。

离散采用食品厚度方向的有限体积。界面潜热记在第一个控制体，不摊进整个域。T2 在截面上增加切向通量与脂肪层的并联热阻。时间格式用半隐式处理扩散、显式处理已冻结的相变速率；步长减半属于 S0，未通过 S0 的网格不进入 S1。

\section{无量纲结构与尺度分离}

取食品厚度 $L$、外部温差 $\Delta T$、初始液相分数 $\varepsilon_{l0}$。Fourier 数 $\mathrm{Fo}=k_{\mathrm{eff}}t/((\rho c)_{\mathrm{eff}}L^2)$ 衡量内部导热是否已经追上边界。Biot 数 $\mathrm{Bi}=hL/k_{\mathrm{eff}}$ 衡量外部对流是否主导表面梯度。Kossovich 数 $\mathrm{Ko}=\Delta h_{\mathrm{vap}}\varepsilon_{l0}\rho_l/((\rho c)_{\mathrm{eff}}\Delta T)$ 衡量蒸发潜热相对显热的大小。蒸发数取界面相变率与内部扩散率之比。

蒸制阶段 $\mathrm{Ko}$ 大且界面冷凝，表面温度先跃迁，冷点仍由 $\mathrm{Fo}$ 控制。短探针的持续时间按表面节点时间常数选取，使其落在冷点 $\mathrm{Fo}\ll 1$ 的窗口内。这个窗口是第 \ref{ch:ati} 章损伤预算的物理来源。

完整边界分类、一维径向离散和两节点与径向模型在 $\vect{z}$ 上的对照，与参数单位一起放在附录。性能数字来自第 \ref{ch:sim} 章，不由无量纲估计代替。
